\chapter{Reti convoluzionali}
\label{cap:cnn}

\begin{inbreve}
Le reti convoluzionali (CNN) mostrano come si progetta una rete feedforward \emph{non} completamente connessa per
un problema specifico, inserendo nell'architettura conoscenza a priori sui dati (le immagini). Il capitolo
presenta la convoluzione come unità che scorre sull'input con pesi condivisi, il pooling, l'architettura
complessiva con i suoi vantaggi, e i risultati che hanno fatto delle CNN il primo esempio di rete profonda.
\end{inbreve}

\section{Il problema e l'idea}

Nel riconoscimento di caratteri scritti a mano (ad esempio immagini $16\times16$ in scala di grigi) un primo approccio
è un MLP standard con 256 input, uno per pixel: si ottiene un errore del 5--20\%, con molti errori dovuti alla
\textbf{mancanza di invarianza} a traslazioni, rotazioni e deformazioni. Una rete completamente connessa tratta ogni
pixel come un input indipendente: non ``sa'' che pixel vicini sono correlati né che uno stesso tratto può comparire
in posizioni diverse.

\begin{idea}[Conoscenza a priori nell'architettura]
\begin{itemize}
  \item \textbf{Connessioni locali} (campi recettivi locali): ogni unità vede solo una piccola zona dell'immagine ed
        estrae \textbf{feature locali}.
  \item \textbf{Condivisione dei pesi} (\ingl{weight sharing}) tra unità diverse: la \textbf{stessa operazione} è
        applicata a parti diverse dell'input, perché le feature di un carattere possono comparire ovunque
        nell'immagine. Si riduce il numero di parametri liberi pur mantenendo una connettività complessa.
\end{itemize}
\end{idea}

\section{La convoluzione}

Il nome viene dall'\termine{operatore di convoluzione}, una media pesata di una funzione $f$ con pesi dati da
un'altra funzione $g$ che scorre:
\[
(f*g)(t) = \int_{-\infty}^{\infty}f(\tau)\,g(t-\tau)\,d\tau .
\]
$g$ è non nulla solo su un intervallo limitato e agisce come una \textbf{finestra scorrevole}.

\subsection{Convoluzione 1D}
\label{sec:conv1d}

Un'unità con tre pesi che scorre su una sequenza di input calcola
\[
\out_t = \sum_{i=1}^{3}w_i\,x_{t+i-2};
\]
ad esempio $o_2 = w_1x_1+w_2x_2+w_3x_3$ e $o_3 = w_1x_2+w_2x_3+w_3x_4$ (\figrif{fig:conv1d}). I pesi sono
\textbf{condivisi} tra le copie dell'unità e si apprendono come al solito: l'unità è un \textbf{filtro adattivo}, e
più unità diverse formano un \textbf{banco di filtri}. Questa rete si chiama anche \ingl{Time-Delay Neural Network}.

\begin{figure}[ht]
\centering
\begin{tikzpicture}[x=1.2cm, y=1cm]
\foreach \i in {1,...,5} \node[neur] (x\i) at (\i,0) {$x_\i$};
\foreach \t in {2,3,4} \node[neurp] (o\t) at (\t,1.6) {$o_\t$};
\foreach \t/\c in {2/blu,3/rosso,4/verde} {
  \pgfmathtruncatemacro{\a}{\t-1} \pgfmathtruncatemacro{\b}{\t+1}
  \draw[->, \c, thick] (x\a) -- (o\t); \draw[->, \c, thick] (x\t) -- (o\t); \draw[->, \c, thick] (x\b) -- (o\t);
}
\node[font=\scriptsize, blu] at (1.25,0.85) {$w_1$}; \node[font=\scriptsize, blu] at (1.82,0.85) {$w_2$}; \node[font=\scriptsize, blu] at (2.75,0.55) {$w_3$};
\node[font=\footnotesize, align=left, anchor=west] at (5.6,0.8) {stessa unità, stessi pesi\\$(w_1,w_2,w_3)$, spostata\\di un passo alla volta};
\end{tikzpicture}
\caption{Convoluzione 1D: la stessa unità (stessi pesi) scorre sulla sequenza di input.}
\label{fig:conv1d}
\end{figure}

\subsection{Convoluzione 2D}

Su un'immagine si usa un \termine{kernel} 2D, ad esempio $3\times3$ (il \textbf{campo recettivo locale}
dell'unità), che si sposta di un pixel alla volta (\termine{stride} 1); il \termine{padding} gestisce i bordi. Il
kernel scorre sull'immagine e produce una \termine{feature map}, le feature estratte da quel filtro. In termini di reti
neurali: si percorre l'immagine con \textbf{lo stesso neurone}, che usa connessioni locali (\figrif{fig:conv2d}).

\begin{figure}[ht]
\centering
\begin{tikzpicture}[x=0.75cm, y=0.75cm]
\foreach \r/\riga in {0/{a,b,c,d}, 1/{e,f,g,h}, 2/{i,j,k,l}} {
  \foreach \v [count=\c from 0] in \riga {
    \draw (\c,-\r) rectangle ++(1,1); \node[font=\small] at (\c+0.5,-\r+0.5) {$\v$};
  }
}
\draw[ultra thick, blu] (0,-1) rectangle (2,1);
\node[font=\footnotesize] at (2,-2.6) {input $3\times4$};
\begin{scope}[xshift=5.2cm]
\foreach \v [count=\c from 0] in {w,x} { \draw (\c,0) rectangle ++(1,1); \node[font=\small] at (\c+0.5,0.5) {$\v$}; }
\foreach \v [count=\c from 0] in {y,z} { \draw (\c,-1) rectangle ++(1,1); \node[font=\small] at (\c+0.5,-0.5) {$\v$}; }
\node[font=\footnotesize] at (1,-2.6) {kernel $2\times2$};
\end{scope}
\begin{scope}[xshift=8.4cm]
\node[draw, fill=blu!10, minimum width=2.6cm, minimum height=0.75cm, font=\scriptsize] at (1.7,0.5) {$aw+bx+ey+fz$};
\node[draw, minimum width=2.6cm, minimum height=0.75cm, font=\scriptsize] at (5.2,0.5) {$bw+cx+fy+gz$};
\node[draw, minimum width=2.6cm, minimum height=0.75cm, font=\scriptsize] at (8.7,0.5) {$cw+dx+gy+hz$};
\node[draw, minimum width=2.6cm, minimum height=0.75cm, font=\scriptsize] at (1.7,-0.5) {$ew+fx+iy+jz$};
\node[draw, minimum width=2.6cm, minimum height=0.75cm, font=\scriptsize] at (5.2,-0.5) {$fw+gx+jy+kz$};
\node[draw, minimum width=2.6cm, minimum height=0.75cm, font=\scriptsize] at (8.7,-0.5) {$gw+hx+ky+lz$};
\node[font=\footnotesize] at (5.2,-2.6) {feature map $2\times3$};
\end{scope}
\end{tikzpicture}
\caption{Convoluzione 2D di un input $3\times4$ con un kernel $2\times2$ (stride 1, senza padding): ogni elemento
della feature map è la somma dei prodotti del kernel con una porzione $2\times2$ dell'input; in blu la prima
posizione.}
\label{fig:conv2d}
\end{figure}

\begin{attenzione}[Convoluzione o cross-correlazione]
La convoluzione 2D di un'immagine $I$ con un kernel $K$ è $S(i,j) = (I*K)(i,j) = \sum_m\sum_nI(m,n)\,K(i-m,j-n)$.
Molte librerie implementano in realtà la \textbf{cross-correlazione}, $S(i,j) = \sum_m\sum_nI(i+m,j+n)\,K(m,n)$,
come nella \figrif{fig:conv2d} (i dettagli non sono richiesti).
\end{attenzione}

Con \textbf{stride} maggiore di 1 il kernel ``salta'' pixel (ad esempio due alla volta): è un
\textbf{sotto-campionamento} e produce una feature map più piccola.

\subsection{Una rete di filtri}

I pesi di un'unità sono un \textbf{filtro addestrato} a rilevare certe feature dell'immagine. Il progetto
dell'architettura prevede:
\begin{itemize}
  \item un \textbf{campo recettivo piccolo e locale}: i filtri appresi rispondono al massimo a pattern
        \textbf{spazialmente locali};
  \item \textbf{lo stesso filtro su tutta l'immagine}: le feature sono rilevate indipendentemente dalla posizione, e
        si ottiene l'\textbf{invarianza per traslazione};
  \item \textbf{filtri appresi}: i filtri che nell'elaborazione delle immagini tradizionale erano progettati a mano
        vengono appresi, senza dipendere da conoscenza a priori e sforzo umano;
  \item \textbf{molti strati impilati}: le unità delle feature map rappresentano aree via via più grandi
        dell'immagine originale (\figrif{fig:campo-recettivo}).
\end{itemize}

\begin{figure}[ht]
\centering
\begin{tikzpicture}[x=0.75cm, y=1.05cm]
\foreach \l in {0,1,2,3} {
  \pgfmathtruncatemacro{\n}{11-2*\l}
  \foreach \i in {1,...,\n} \node[circle, draw, inner sep=0pt, minimum size=4mm] (n\l-\i) at (\i+\l,\l) {};
}
\fill[rosso] (n3-3) circle (2.2mm);
\foreach \i in {3,4,5} \draw[rosso, thick] (n3-3) -- (n2-\i);
\foreach \i in {3,4,5} \fill[rosso!55] (n2-\i) circle (1.6mm);
\foreach \i in {3,...,7} \fill[rosso!35] (n1-\i) circle (1.6mm);
\foreach \i in {3,...,9} \fill[rosso!20] (n0-\i) circle (1.6mm);
\draw[rosso, densely dashed] (n3-3) -- (n0-3.south west);
\draw[rosso, densely dashed] (n3-3) -- (n0-9.south east);
\node[font=\footnotesize, anchor=west] at (12.2,0) {input};
\node[font=\footnotesize, anchor=west] at (12.2,1) {strato 1};
\node[font=\footnotesize, anchor=west] at (12.2,2) {strato 2};
\node[font=\footnotesize, anchor=west] at (12.2,3) {strato 3};
\end{tikzpicture}
\caption{Crescita del campo recettivo con la profondità (caso 1D, kernel di ampiezza 3): un'unità dello strato 3 vede
3 unità dello strato 2 ma, indirettamente, 7 input.}
\label{fig:campo-recettivo}
\end{figure}

\section{Pooling}

Il \termine{pooling} riduce la feature map producendo un valore per un insieme rettangolare di posizioni invece che
per ogni posizione: con un sotto-campionamento (stride $>1$), con una \textbf{media} (semplice o pesata) o, più
comunemente, con il \termine{max pooling} (\figrif{fig:maxpool}). Il pooling rende la rappresentazione
\textbf{approssimativamente invariante a piccole traslazioni}: il massimo (o la media) di una zona cambia poco se
l'input si sposta di poco.

\begin{figure}[ht]
\centering
\begin{tikzpicture}[x=0.7cm, y=0.7cm]
\foreach \r/\riga in {0/{1,1,2,4}, 1/{5,6,7,8}, 2/{3,2,1,0}, 3/{1,2,3,4}} {
  \foreach \v [count=\c from 0] in \riga {
    \pgfmathtruncatemacro{\q}{(\c>1)+2*(\r>1)}
    \ifcase\q \def\col{blu!15}\or \def\col{rosso!15}\or \def\col{verde!18}\or \def\col{arancio!20}\fi
    \fill[\col] (\c,-\r) rectangle ++(1,1); \draw (\c,-\r) rectangle ++(1,1);
    \node[font=\small] at (\c+0.5,-\r+0.5) {\v};
  }
}
\draw[->, thick] (4.4,-1) -- (6.2,-1) node[midway, above, font=\scriptsize, align=center] {max pooling\\$2\times2$, stride 2};
\foreach \v/\c/\r/\col in {6/0/0/blu!15, 8/1/0/rosso!15, 3/0/1/verde!18, 4/1/1/arancio!20} {
  \fill[\col] (6.8+\c,-0.5-\r) rectangle ++(1,1); \draw (6.8+\c,-0.5-\r) rectangle ++(1,1);
  \node[font=\small] at (7.3+\c,-\r) {\v};
}
\end{tikzpicture}
\caption{Max pooling con filtro $2\times2$ e stride 2: ogni blocco viene sostituito dal suo massimo.}
\label{fig:maxpool}
\end{figure}

\section{La CNN nel complesso}

Una CNN realizza con la condivisione dei pesi una finestra scorrevole di campi locali sul segnale, estesa alle
immagini 2D e ripetuta su molti strati di feature map, alternata a operazioni di pooling; in cima ci sono di solito
alcuni strati completamente connessi (con softmax per la classificazione). Il numero di filtri può crescere negli
strati più alti; per un'immagine a colori (RGB) l'input ha tre canali, e la ``profondità'' del volume di ogni strato è
il suo numero di feature map (\figrif{fig:cnn}).

\begin{figure}[ht]
\centering
\begin{tikzpicture}[x=1cm, y=1cm]
\newcommand{\pila}[5]{% x, larghezza, altezza, numero, etichetta
  \foreach \k in {1,...,#4} \draw[fill=white] (#1+0.12*\k, 0.12*\k) rectangle ++(#2,#3);
  \node[font=\scriptsize, align=center, anchor=north] at (#1+0.06*#4+#2/2, -0.1) {#5};}
\pila{0}{1.6}{1.6}{1}{input}
\pila{2.3}{1.3}{1.3}{4}{convoluzione\\(+ReLU)}
\pila{4.6}{0.75}{0.75}{4}{pooling}
\pila{6.4}{0.6}{0.6}{7}{convoluzione}
\pila{8.4}{0.32}{0.32}{7}{pooling}
\foreach \y in {0.2,0.5,...,1.4} \fill (10.2,\y) circle (1.5pt);
\node[font=\scriptsize, anchor=north, align=center] at (10.2,-0.1) {strati\\densi};
\foreach \y in {0.5,0.8,1.1} \draw (11.2,\y) circle (2.5pt);
\node[font=\scriptsize, anchor=north, align=center] at (11.2,-0.1) {uscita\\(softmax)};
\draw[decorate, decoration={brace, amplitude=5pt}] (2.3,2.4) -- (9.4,2.4) node[midway, above=5pt, font=\footnotesize] {feature learning};
\draw[decorate, decoration={brace, amplitude=5pt}] (9.9,2.4) -- (11.5,2.4) node[midway, above=5pt, font=\footnotesize] {classificazione};
\end{tikzpicture}
\caption{Schema di una CNN: convoluzioni e pooling alternati (estrazione delle feature) seguiti da strati completamente
connessi (classificazione). Le feature map si rimpiccioliscono e aumentano di numero salendo.}
\label{fig:cnn}
\end{figure}

\begin{idea}[I quattro ingredienti e i loro vantaggi]
\begin{enumerate}
  \item \textbf{Connessioni locali} e
  \item \textbf{pesi condivisi}: rilevano pattern locali in modo \textbf{invariante alla traslazione} e
        \textbf{riducono il numero di parametri liberi} pur mantenendo le connessioni, il che è una forma di
        \textbf{regolarizzazione}.
  \item \textbf{Pooling}: riduce la dimensione della rappresentazione a ogni strato (una ``piramide'' di strati) e
        aiuta l'invarianza a piccoli spostamenti e distorsioni.
  \item \textbf{Molti strati}: le operazioni 1 e 2 applicate a tutte le feature map producono un'\textbf{astrazione
        progressiva}, che compone primitive di basso livello (apprese una volta sola) nelle combinazioni utili agli
        strati successivi. Le unità degli ultimi strati hanno un campo recettivo che copre gran parte dell'immagine.
\end{enumerate}
\end{idea}

L'addestramento avviene con la \textbf{backpropagation}, con le euristiche note e qualche specializzazione (per i pesi
condivisi, il pooling, \dots). Vista la dimensione tipica delle reti e dei dati, molti iperparametri si fissano
\textbf{per esperienza}, perché una cross-validation su grandi spazi di iperparametri sarebbe troppo costosa.

\section{Risultati}

\subsection{Pesi condivisi e connessioni locali}

Su un problema di cifre manoscritte (circa 300 immagini di training) LeCun confrontò cinque reti: Net-1 senza strati
nascosti, Net-2 completamente connessa con uno strato nascosto, Net-3 con connessioni locali $3\times3$, Net-4 e Net-5
con connessioni locali e pesi condivisi.

\begin{table}[ht]
\centering\small
\caption{Connessioni e pesi delle cinque reti.}
\begin{tabular}{lrr}
\toprule
Rete & Connessioni & Pesi \\
\midrule
Net-1 & 2570 & 2570 \\
Net-2 & 3214 & 3214 \\
Net-3 & 1226 & 1226 \\
Net-4 & 2266 & 1132 \\
Net-5 & 5194 & 1060 \\
\bottomrule
\end{tabular}
\end{table}

Con i pesi condivisi il numero di \textbf{connessioni} può crescere mentre il numero di \textbf{pesi} (parametri liberi)
resta basso. L'errore di training è nullo per tutte le reti, mentre l'accuratezza di test va da circa il 72\% (Net-1)
e l'88\% (Net-2, Net-3) fino al 98\% (Net-4) e quasi al 99\% (Net-5, errore dell'1,6\%).

\subsection{Altri risultati}

\begin{itemize}
  \item Sul dataset MNIST (60\,000 immagini di training, 10\,000 di test) le reti profonde, convoluzionali e no,
        sono arrivate a errori di pochi decimi di punto percentuale. Anche un semplice MLP con molti strati e molte
        unità ottiene risultati di questo livello se si controlla la complessità generando \textbf{molte immagini
        deformate} per il training (\ingl{data augmentation}).
  \item AlexNet, una CNN profonda (cinque strati convoluzionali, alcuni seguiti da max pooling, e strati
        completamente connessi con softmax a 1000 classi, circa 60 milioni di parametri), migliorò nettamente lo
        stato dell'arte nella classificazione di immagini (ImageNet). Per velocizzare l'addestramento usava unità
        \textbf{non saturanti} (ReLU) e le GPU; per ridurre l'overfitting negli strati densi usava il
        \textbf{dropout} (\rif{sec:dropout}).
\end{itemize}

\begin{attenzione}[Un test set usato troppe volte]
Il test set di MNIST è stato usato così tante volte dalla comunità che miglioramenti dell'ordine dello 0,1\% non sono
più significativi: di fatto è diventato un validation set (\rif{sec:errori-validazione}).
\end{attenzione}

Esistono molte architetture specializzate per le immagini, e si possono riusare le feature di reti
\textbf{pre-addestrate} su grandi dataset (\rif{sec:transfer}). Le implementazioni efficienti sfruttano le
\textbf{GPU} tramite la rappresentazione a \textbf{tensori}: la moltiplicazione di matrici (cioè i prodotti scalari) si
parallelizza molto bene, e molte operazioni di convoluzione si possono scrivere come un'unica moltiplicazione tra
matrici a blocchi o con operazioni su tensori.

Un'alternativa al max pooling sono le \textbf{Capsule Network}: una ``capsula'' è un gruppo di neuroni che produce sia
un parametro che indica se un'entità è presente, sia un vettore di \textbf{parametri di posa} (posizione,
orientamento); al posto del pooling un meccanismo di \textbf{routing} instrada l'informazione di basso livello verso la
capsula più adatta.

\begin{puntichiave}
\begin{itemize}
  \item Convoluzione = la stessa unità (filtro adattivo) che scorre sull'input: connessioni locali e pesi condivisi.
  \item Kernel, stride, padding, feature map; il pooling (soprattutto max) riduce la mappa e dà invarianza a piccoli
        spostamenti.
  \item I pesi condivisi riducono i parametri liberi: è una regolarizzazione.
  \item Con la profondità il campo recettivo cresce e le feature diventano più astratte.
\end{itemize}
\end{puntichiave}

\begin{domande}
\begin{enumerate}
  \item Quali sono le idee chiave delle CNN e cosa comporta ciascuna?
  \item Cos'è la convoluzione in una rete neurale? Cosa sono kernel, stride, padding, feature map?
  \item A cosa serve il pooling? Come funziona il max pooling?
  \item Perché la condivisione dei pesi è una forma di regolarizzazione?
  \item Come cresce il campo recettivo con la profondità?
\end{enumerate}
\end{domande}
